% !PRINTED 525
\weberpage{525}
Now let $\pi_1, \pi_2, \pi_3, \ldots$ be an arbitrary number of given prime functionals, different from one another. We set

\[
\omega_1 = \omega\pi_2\pi_3\ldots, \quad \omega_2 = \omega\pi_1\pi_3\ldots, \quad \omega_3 = \omega\pi_1\pi_2\ldots,
\]
and determine, by what has just been proved, the integers $\alpha_1, \alpha_2, \alpha_3, \ldots$ in $\Omega$, so that
\[
\begin{array}{l}
\alpha_1 \text{ becomes divisible by } \omega_1 \text{, but not by } \omega_1\pi_1\\
\alpha_2 \quad \text{\textquotedbl} \quad \text{\textquotedbl} \quad \text{\textquotedbl} \quad \omega_2 \quad \text{\textquotedbl} \quad \text{\textquotedbl} \quad \text{\textquotedbl} \quad \omega_2\pi_2\\
\alpha_3 \quad \text{\textquotedbl} \quad \text{\textquotedbl} \quad \text{\textquotedbl} \quad \omega_3 \quad \text{\textquotedbl} \quad \text{\textquotedbl} \quad \text{\textquotedbl} \quad \omega_3\pi_3\\
\cdots\cdots\cdots\cdots\cdots\cdots\cdots\cdots\cdots\cdots\cdots\cdots\cdots\cdots\cdots\cdots,
\end{array}
\]
and from these derive the integer
\[
\alpha = \alpha_1 + \alpha_2 + \alpha_3 + \cdots
\]
This number is evidently divisible by $\omega$, since all the summands $\alpha_1, \alpha_2, \alpha_3, \ldots$ are divisible by $\omega$. But it is not divisible by $\omega\pi_1$, because $\alpha_2, \alpha_3, \ldots$, though not $\alpha_1$, are divisible by $\omega\pi_1$; and likewise it is not divisible by $\omega\pi_2, \omega\pi_3, \ldots$. If, therefore, we set
\[
\alpha = \omega\eta
\]
then $\eta$ is an integral functional which is not divisible by $\pi_1, \pi_2, \pi_3, \ldots$, and which therefore, when $\pi_1, \pi_2, \pi_3, \ldots$ are the prime factors of $\mu$ different from one another, is relatively prime to $\mu$.

Now let us assume, arbitrarily, an integer $\beta$ in $\Omega$ divisible by $\omega$, and let us set $\beta = \omega\mu$; then we can determine $\alpha = \omega\eta$ so that $\eta$ becomes relatively prime to $\mu$, and then $\omega$ is the greatest common divisor of $\alpha$ and $\beta$. From this it follows:

4. \gesperrt{Every integral functional $\omega$ of the field $\Omega$ is the greatest common divisor of two integers, and consequently $\omega$ is associated with a binary linear form $\alpha x + \beta y$, in which $\alpha$ and $\beta$ are integers.}

\begin{center}
\rule{4cm}{0.4pt}
\end{center}

A general remark, one that reaches over into another region, may still find its place here.

It is the principal result of this chapter that the integral algebraic numbers in a definite field can be decomposed into prime factors in a unique manner, exactly in the same way as is known for the integral rational numbers; though indeed only by the fact that the content of the field
%CONT>>
